\documentstyle[% 


righttag% 


,11pt% 


,amssymb% 


,russian%               remove in Eng. 


,izvru%                 Set izven% in Eng. 


]{amsart} 


 


% attn to \wt vs \Tilde


 


\newcommand{\re}{\operatorname{Re}} 


\newcommand{\im}{\operatorname{Im}} 


\newcommand{\Res}{\operatornamewithlimits{res}} 


\newcommand{\tts}[1]{} 


 


\theoremstyle{plain} 


\theorembodyfont{\it} 


\newtheorem{thmnonum}{\hskip\parindent Теорема} 


\newtheorem{lemnonum}{\hskip\parindent Лемма} 


\renewcommand{\thethmnonum}{} 


\renewcommand{\thelemnonum}{} 


 


%\hoffset=-15mm%                  For Eng. 


\setcounter{page}{48} 


\god{2002} 


\nomer{10~(485)} %                        Remove in Eng. 


%\nomer{Vol.\,46, No.\,10}%               Set in Eng. 


%\str{\pageref{begin}--\pageref{end}}%  Set in Eng. 


\udk{517.544} 


 


\author{\it Г.Р.\,Галиуллина, С.Р.\,Насыров} 


% NB:   ^^ remove in Eng. 


\title{Уравнение Гахова для внешней смешанной 


обратной краевой задачи по параметру $\boldmath x$} 


\thanks{Работа второго автора поддержана грантом Российского фонда 


фундаментальных исследований\linebreak \No\,02-01-00914.} 


 


\date{} 


%\pagestyle{myheadings}%         Set in Eng. 


\pagestyle{plain} %              Remove in Eng. 


\begin{document} 


%\thispagestyle{plain}%          Set in Eng. 


\maketitle 


\label{begin} 


 


 


\section*{Введение} 


 


Внутренняя задача, соответствующая рассматриваемой в данной статье, была 


впервые 


поставлена В.Н.\,Монаховым (см. [1]). В [2] была доказана разрешимость 


внутренней задачи на римановых поверхностях, исследован 


характер неединственности решения в зависимости от геометрии известной 


части границы (по поводу истории вопроса см. подробнее [1], [2]). 


 


В данной работе рассматривается случай, когда неизвестная область 


$D_z$ в плоскости $z$ содержит бесконечно удаленную точку. В случае, 


когда 


известная часть границы искомой области полигональна, получено 


интегральное 


представление решения $z(\zeta)$, $\zeta\in D_\zeta$, где $D_\zeta$ --- 


верхняя полуплоскость 


в $\zeta$-плоскости. Как и во внутренней задаче, при доказательстве 


разрешимости возникает проблема определения акцессорных параметров, 


но здесь эта трудная и интересная задача не исследуется. Основное 


внимание уделяется однозначности функции $z(\zeta)$, которая 


определяется 


полученным интегральным представлением. Необходимым и достаточным 


условием этой однозначности является равенство нулю вычета производной 


функции $dz(\zeta)/d\zeta$ в точке $\zeta_0$, которая соответствует 


бесконечно 


удаленной точке в плоскости $z$. 


 


Это равенство можно рассматривать как уравнение относительно 


переменной $\zeta_0$. Во внешней обратной краевой задаче по параметру 


$s$ соответствующее уравнение впервые получил Ф.Д.\,Гахов [4], который 


доказал его разрешимость. В дальнейшем уравнение Гахова и его обобщения 


рассматривались многими авторами (см., напр., [3]--[5]), которые 


изучали 


вопросы единственности решения этого уравнения, структуру множества 


его корней и пр. В частности, в [6] показано, что уравнение Гахова во 


внешней обратной краевой задаче по параметру $s$ имеет конечное число 


решений. 


 


Основным результатом статьи является доказательство разрешимости 


аналога уравнения Гахова для внешней смешанной обратной краевой 


задачи по параметру $x$ в случае полигональной известной части границы. 


Кроме того, с использованием результатов из теории полианалитических 


функций исследована локальная структура множества решений 


этого уравнения. 


 


 


\section{Постановка задачи} 


 


Пусть $D_z$ --- односвязная жорданова область на сфере Римана, 


содержащая 


внутри себя $\infty$, с границей $L_z$, которая состоит из известной 


дуги 


$L^1_z$ и искомой дуги $L^2_z$. В дальнейшем будем считать, что $L^1_z$ 


--- полигон с 


вершинами $z_k=x_k+i y_k$, $k=1,\dotsc,n$. Предполагается, что $L^2_z$ 


такова, 


что любая прямая, параллельная мнимой оси, пересекает ее не более, чем 


в одной точке, и $x_1=a<b=x_n$ (рис.~1). 


 


Требуется найти $L^2_z$ и аналитическую в области $D_z$ функцию $w(z)$, 


конформно отображающую $D_z$ на жорданову область $D_w$ и 


удовлетворяющую 


следующим краевым условиям. 


 


B плоскости $w=\varphi+i\psi$ дуге $L^2_z$ соответствует дуга $L^2_w$ с 


уравнением $\varphi=f_1(x)$, $\psi=f_2(x)$, где $w(x)=f_1(x)+if_2(x)$, 


$x\in[a,b]$, --- граничные значения искомой аналитической 


функции $w(z)$ на $L^2_z$ и $x=\re z$. Будем предполагать, что 


функция $w(x)$ непрерывно дифференцируема и $w'(x)\ne0$, $a\le x\le b$. 


Уравнение дуги $L^1_w$, дополняющей $L^2_w$ до замкнутого контура 


$L_w$, $\Phi(\varphi,\psi)=0$ 


считается заданным. Предполагается, что функция $\Phi(\varphi,\psi)$ 


дважды непрерывно 


дифференцируема и гладкие дуги $L^1_w$ и $L^2_w$ образуют в точках 


стыка $w_1$ и $w_n$ ненулевые углы $\pi\gamma_1$ и $\pi\gamma_n$. 





\begin{center} 


\begin{picture}(470,320)(0,-15) 


\put(85,290){\special{em:graph gal-nas.bmp}} 


\put(212,-10){\mbox Рис.\,1.} 


\end{picture} 


\end{center} 


 


Конформно отобразим верхнюю полуплоскость $D_\zeta=\{\im\zeta>0\}$ на 


$D_w$ функцией $w=\omega(\zeta)$ так, чтобы точки $\infty$, $t_1=-1$, 


$t_n=1$, лежащие 


на вещественной оси, переходили соответственно в фиксированную 


точку $w_0\in L^1_w$, $w_1$ и $w_n$. Пусть $t_k$ --- точки на границе 


$D_\zeta$, соответствующие 


вершинам ломаной $L^1_z$, $k=1,\dotsc,n$. Сравнивая граничные значения 


функций $w(x)=f_1(x)+if_2(x)$ и $w=\omega(t)$ на участках, 


соответствующих $L^2_w$, 


получим соотношение $f_1(x)+if_2(x)=\omega(t)$, из которого найдем 


зависимость $x=H(t)$. Дифференцируя ее по $t$, получим


\begin{equation} 


\frac{dx(t)}{dt}=\frac{\omega'(t)}{w'(H(t))}=h(t),\quad |t|<1, 


\end{equation} 


где $h(t)=H'(t)\le0$ при $t\in(-1,1)$, т.\,к. функция $x(t)$ монотонно 


убывает. Вследствие достаточной гладкости дуги $L^2_w$ функция 


$d\omega/dt$ на 


$(-1,1)$ особенностей не имеет, т.\,е. $0\ne|d\omega/dt|<\infty$ при 


$|t|<1$. Таким 


образом, $-\infty<h(t)<0$, $|t|<1$. 


 


Учитывая поведение производной отображающей функции $d\omega/dt$ в 


угловых 


точках $t_1$ и $t_n$ (см., напр., [1]), представим $h(t)$ в виде 


$h(t)=h^*(t)(t-t_1)^{\gamma_1-1}(t-t_n)^{\gamma_n-1}$, 


где $h^*(t)$ --- г\"ельдерова функция на $[-1,1]$, $-\infty<h^*(t)<0$, 


$t\in[-1,1]$. 


 


Пусть уравнения прямых, на которых лежат стороны полигона, 


имеют вид $a_k x-b_k y=c_k$, $k=1,\dotsc,n-1$. 


Тогда 


\begin{equation} 


a_k x(t)-b_k y(t)=c_k,\quad t\in(t_k,t_{k+1}),\ \ k=1,\dotsc,n-1. 


\end{equation} 


Отметим, что здесь через $(t_k,t_{k+1})$ обозначена часть границы 


$D_\zeta$ в 


расширенной комплексной плоскости от точки $t_k$ до точки $t_{k+1}$, 


проходимая 


в положительном направлении. 


 


Пусть $\zeta_0$ --- точка в $D_\zeta$, соответствующая точке $\infty$ в 


плоскости $z$. 


Найдем функцию $z=F(\zeta)$, отображающую верхнюю полуплоскость 


$D_\zeta$ на 


область $D_z$, которая на вещественной оси удовлетворяет краевым 


условиям (1), (2). 


Если она известна, то известно и уравнение $z=F(t)$, $-1\le t\le1$ 


контура $L^2_z$, а функция $w(z)$ восстанавливается по формуле 


$w=\omega(F^{-1}(z))$, 


где $\zeta=F^{-1}(z)$ --- функция, обратная к $z=F(\zeta)$. 


Проведем из точек $z_1$ и $z_n$ вверх вертикальные лучи $l_1$ и $l_n$ и 


обозначим 


через $\pi\alpha_1$, $\pi\alpha_n$ углы, образованные первым и 


$(n-1)$-м звеньями ломаной 


$L^1_z$ с этими лучами. Пусть $\pi\alpha_k$ --- внутренние углы области 


$D_z$ в точках 


$z_k$, $k=2,\dotsc,n-1$ (рис.~1). 


 


По аналогии с внутренней задачей (см. [2]) нетрудно показать, что 


необходимым условием разрешимости задачи является следующее ограничение 


на известную часть границы: ломаная $L^1_z$, дополненная лучами 


$l_1$ и $l_n$, является границей двулистной многоугольной римановой 


поверхности без точек ветвления, причем 


$$ 


\sum^n_{k=1}(\alpha_k-1)=3. 


$$ 


В дальнейшем будем предполагать это условие выполненным. 


 


\section{Построение интегрального представления решения} 


 


Дифференцируя, запишем (1), (2) в виде


\begin{gather*} 


a_k\frac{dx(t)}{dt}-b_k\frac{dy(t)}{dt}=0,\quad 


t\in[t_k,t_{k+1}],\ \ k=1,\dotsc,n-1, \\ 


% 


\frac{dx(t)}{dt}=h(t),\quad t\in[-1,1]. 


\end{gather*} 


Эта задача является краевой задачей Гильберта с разрывными 


коэффициентами 


для функции $dz(\zeta)/d\zeta$ в полуплоскости 


$$ 


\re[(a(t)+i b(t))dz(t)/dt]=c(t),\quad t\in(-\infty,\infty), 


$$ 


где 


\begin{align*} 


a(t)&= 


\begin{cases} 


a_k, & t\in[t_k,t_{k+1}],\ \; k=\ov{1,n-1};\\ 


1, & t\in[-1,1], 


\end{cases} 


\\ 


b(t)&= 


\begin{cases} 


b_k, & t\in[t_k,t_{k+1}],\ \; k=\ov{1,n-1};\\ 


0, & t\in[-1,1], 


\end{cases} 


\\ 


c(t)&= 


\begin{cases} 


0, & |t|>1;\\ 


h(t), & t\in[-1,1]. 


\end{cases} 


\end{align*} 


 


Решение ищем в классе функций $dz(\zeta)/d\zeta$, ограниченных в 


вершинах 


полигона с углами, б\'ольшими $\pi$, неограниченных в остальных и 


имеющих полюс второго порядка в точке $\zeta_0$. 


 


Перепишем задачу Гильберта в виде $\re[(a(t)+i 


b(t))\Phi(t)]=c(t)|t-\zeta_0|^4$, 


гдe 


$\Phi(\zeta)=(\zeta-\zeta_0)^2(\zeta-\ov\zeta_0)^2dz(\zeta)/d\zeta$. 


 


Пусть $\alpha_1<1$, $\alpha_n<1$. Функция 


$\Pi(\zeta)=\prod\limits^n_{k=1}(\zeta-t_k)^{\alpha_k-1}$ может быть 


взята за каноническую функцию однородной задачи, и решение запишется 


так


$$ 


\Phi(\zeta)=\frac{dz(\zeta)}{d\zeta} 


(\zeta-\zeta_0)^2(\zeta-\ov\zeta_0)^2=\frac{\Pi(\zeta)}{\pi i} 


\int^1_{-1}\frac{h(t)|t-\zeta_0|^4}{\Pi(t)(t-\zeta)}dt 


+P(\zeta)\Pi(\zeta), 


$$ 


где $P(\zeta)$ --- некоторый многочлен. Нетрудно видеть, что на 


бесконечности $|\Pi(\zeta)|\sim|\zeta|^d$, где 


$d=\sum\limits^n_{k=1}(\alpha_k-1)=3$. 


 


Чтобы искомый контур был конечен, следует положить $P(\zeta)\equiv0$, 


т.\,к. 


$$ 


\frac{\Pi(\zeta)}{(\zeta-\zeta_0)^2(\zeta-\ov\zeta_0)^2} 


\sim\frac{1}{\zeta},\quad \zeta\to\infty. 


$$ 


Итак, 


\begin{equation} 


z=F(\zeta)=z_1+\frac{1}{\pi i}\int^\zeta_{-1} 


\frac{\Pi(\zeta)d\zeta}{(\zeta-\zeta_0)^2(\zeta-\ov\zeta_0)^2} 


\int^1_{-1}\frac{h(t)|t-\zeta_0|^4}{\Pi(t)(t-\zeta)}dt 


\end{equation} 


--- единственное решение задачи в классе ограниченных функций 


$z=z(\zeta)$, 


если $a_j<1$, $j=1,n$. 


 


Пусть теперь только одно из $\alpha_j$, $j=1,n$, например, $\alpha_n$, 


больше 


единицы. Каноническая функция может быть взята в виде 


$\Pi(\zeta)(\zeta-t_n)^{-1}$, 


и тогда интегральное представление решения имеет вид 


\begin{equation} 


\frac{dz(\zeta)}{d\zeta}=\frac{\Pi(\zeta)} 


{(\zeta-\zeta_0)^2(\zeta-\ov\zeta_0)^2}\bigg[\frac{1}{\pi i} 


\int^1_{-1}\frac{h(t)|t-\zeta_0|^4}{\Pi(t)(t-\zeta)}dt+ 


\frac{i K}{\zeta+1}\bigg], 


\end{equation} 


где $K$ --- вещественное число (можно показать, что $K\ge0$). 


Если же $\alpha_1>1$ и $\alpha_n>1$, то имеем


\begin{equation} 


\frac{dz(\zeta)}{d\zeta}=\frac{\Pi(\zeta)} 


{(\zeta-\zeta_0)^2(\zeta-\ov\zeta_0)^2}\bigg[\frac{1}{\pi i} 


\int^1_{-1}\frac{h(t)|t-\zeta_0|^4}{\Pi(t)(t-\zeta)}dt+ 


\frac{i K_1}{\zeta+1}+\frac{i K_2}{\zeta-1}\bigg], 


\end{equation} 


где $K_1$, $K_2$ вещественны. 


Если положить $K=0$ в (4) или $K_1=K_2=0$ в (5), то получим 


решение задачи, которое по виду совпадает с (3). Таким образом, (3) 


является частным решением задачи и в случаях, когда по крайней 


мере один из параметров $\alpha_1$, $\alpha_n$ больше единицы. 


 


\section{Разрешимость уравнения Гахова} 


 


Условием однозначности функции $z(\zeta)$, определенной формулой (3), 


будет служить равенство $c_{-1}=0$, где $c_{-1}$ --- вычет функции 


$dz(\zeta)/d\zeta$ в 


точке $\zeta=\zeta_0$. Имеем 


\begin{multline*} 


\Res_{\zeta=\zeta_0}\frac{dz(\zeta)}{d\zeta}=\lim_{\zeta\to\zeta_0} 


\frac{d}{d\zeta}\bigg[(\zeta-\zeta_0)^2\frac{dz(\zeta)}{d\zeta}\bigg] 


=\lim_{\zeta\to\zeta_0}\frac{d}{d\zeta}\bigg[\frac{1} 


{(\zeta-\ov\zeta_0)^2}\frac{\Pi(\zeta)}{\pi i} 


\int^1_{-1}\frac{h(t)|t-\zeta_0|^4}{\Pi(t)(t-\zeta)}dt\bigg]= \\ 


% 


=\frac{\Pi(\zeta_0)}{(\zeta_0-\ov\zeta_0)^2}\frac{1}{\pi i}\bigg\{ 


-\frac{2}{\zeta_0-\ov\zeta_0} 


\int^1_{-1}\frac{h(t)|t-\zeta_0|^4}{\Pi(t)(t-\zeta_0)}dt+ \\ 


% 


+\sum^n_{j=1}\frac{\beta_j}{\zeta_0-t_j} 


\int^1_{-1}\frac{h(t)|t-\zeta_0|^4}{\Pi(t)(t-\zeta_0)}dt+ 


\int^1_{-1}\frac{h(t)|t-\zeta_0|^4}{\Pi(t)(t-\zeta_0)}\bigg\}, 


\end{multline*} 


где $\beta_j=\alpha_j-1$, $j=1,\dotsc,n$. 


 


Введем обозначения 


$$ 


M(\zeta_0)= 


\int^1_{-1}\frac{h(t)|t-\zeta_0|^4}{\Pi(t)(t-\zeta_0)}dt,\qquad 


N(\zeta_0)= 


\int^1_{-1}\frac{h(t)|t-\zeta_0|^4}{\Pi(t)(t-\zeta_0)^2}dt. 


$$ 


Приравняем вычет $\Res_{\zeta=\zeta_0}dz(\zeta)/d\zeta$ к нулю 


$$ 


\frac{\Pi(\zeta_0)}{(\zeta-\ov\zeta_0)^2}\bigg\{\bigg[ 


-\frac{2}{\zeta_0-\ov\zeta_0}+\sum^n_{k=1}\frac{\beta_k}{\zeta_0-t_k} 


\bigg]M(\zeta_0)+N(\zeta_0)\bigg\}=0. 


$$ 


 


Ясно, что $\Pi(\zeta)=\prod\limits^n_{k=1}(\zeta-t_k)^{\beta_k}\ne0$ 


при $\zeta\ne t_k$, $k=\ov{1,n}$. 


Покажем, что $M(\zeta_0)\ne0$, $\zeta_0\in D_\zeta$. Имеем 


$$ 


M(\zeta)=\int^1_{-1}\frac{h(t)|t-\zeta_0|^4}{\Pi(t)} 


\frac{dt}{t-\zeta_0}=\int^1_{-1} 


\frac{h(t)|t-\zeta_0|^2(t-\ov\zeta_0)}{\Pi(t)}dt, 


$$ 


следовательно, 


$$ 


\im M(\zeta_0)=\im \zeta_0\int^1_{-1} 


\frac{h(t)|t-\zeta_0|^2}{\Pi(t)}dt\ne0,\quad \zeta_0\in D_\zeta. 


$$ 


Отсюда видно, что $M(\zeta_0)\ne0$, $\zeta_0\in D_\zeta$. Таким 


образом, имеем следующее 


условие однозначности функции $z(\zeta)$: 


\begin{equation} 


F(\zeta_0)=-\frac{N(\zeta_0)}{M(\zeta_0)}+\frac{2}{\zeta_0-\ov\zeta_0} 


-\sum^n_{j=1}\frac{\beta_j}{\zeta_0-t_j}=0. 


\end{equation} 


 


Назовем (6) уравнением Гахова для внешней смешанной обратной 


краевой задачи по параметру $x$. Выясним, существует ли решение данного 


уравнения в верхней полуплоскости. Для этого рассмотрим окружность 


радиуса $R$ ($R$ --- достаточно большое число) с центром в начале 


координат. Проведем хорду $T_1$ параллельно оси $O\zeta$ через точку 


$i\wt\eta_0$, где 


$\wt\eta_0$ --- фиксированное достаточно малое число. (Значения $R$ и 


$\wt\eta_0$ будут 


выбраны ниже.) Хорда $T_1$ делит окружность на две части. Рассмотрим 


ту часть $T_2$, которая лежит целиком в верхней полуплоскости. Пусть 


$Q$ --- 


круговой сегмент, ограниченный $T_1$ и $T_2$. 


 


Левая часть (6) определяет некоторое векторное поле на плоскости. 


Подсчитаем его вращение на $\partial Q=T_1\cup T_2$. Нетрудно 


убедиться, что 


функция $N(\zeta_0)/M(\zeta_0)$ непрерывна по $\zeta_0$ в $D_\zeta$ и 


$N(\zeta_0)\sim a_0\ov\zeta{}^2_0$, $M(\zeta_0)\sim 


-a_0\zeta_0\ov\zeta{}^2_0$ 


при $\zeta_0\to\infty$, где 


$$ 


a_0=\int^1_{-1}\frac{h(t)}{\Pi(t)}dt. 


$$ 


(Так как функции $h$ и $\Pi$ непрерывны, не обращаются в нуль на 


$(-1,1)$ 


и имеют постоянный аргумент на этом интервале, то $a_0\ne0$.) 


Следовательно, $N(\zeta_0)/M(\zeta_0)\sim -1/\ov\zeta_0$, 


$\zeta_0\to\infty$. Значит, существует константа 


$C>0$ такая, что 


$$ 


\bigg|\frac{N(\zeta_0)}{M(\zeta_0)}\bigg|\le C,\quad \zeta_0\in 


D_\zeta. 


$$ 


 


1. Рассмотрим поведение $F(\zeta_0)$ при $\zeta_0\in T_1$. 


 


а) Пусть $\zeta_0=\xi_0+i\wt\eta_0\in T_1$ и для некоторого $k$ 


величина $|\xi_0-t_k|<\varepsilon$, 


где $\varepsilon=\frac{1}{3}\min\limits_{1\le j\le n-1}|t_{j+1}-t_j|$.


Тогда 


$$ 


\bigg|\frac{\beta_k}{\zeta_0-t_k} \bigg|= 


\bigg|\frac{\beta_k}{\xi_0-t_k+i\wt\eta_0} \bigg| 


\le\frac{\wt\beta}{\wt\eta_0}, 


$$ 


где $\wt\beta=\max\limits_j|\beta_j|<1$. Имеем 


\begin{multline*} 


\bigg|\sum_{j\ne k}\frac{\beta_j}{\zeta_0-t_j} \bigg|\le 


\sum_{j\ne k}\bigg|\frac{\beta_j}{\xi_0-t_j+i\wt\eta_0} \bigg| 


\le\sum_{j\ne k}\frac{1}{\sqrt{(\xi_0-t_j)^2+\wt\eta{}^2_0}}= \\ 


% 


=\sum_{j\ne k}\frac{1}{\sqrt{(\xi_0-t_k+t_k-t_j)^2+\wt\eta{}^2_0}} 


\le\sum_{j\ne k}\frac{1}{\sqrt{(|t_k-t_j|-|\xi_0-t_k|)^2+ 


\wt\eta{}^2_0}} \le \\ 


% 


\le\sum_{j\ne k}\frac{1}{\sqrt{\varepsilon^2(3|k-j|-1)^2 


+\wt\eta^2_0}}<\frac{n-1}{\sqrt{2}\varepsilon}. 


\end{multline*} 


 


Введем обозначение 


$$ 


F_1(\zeta_0)=\frac{N(\zeta_0)}{M(\zeta_0)}- 


\sum^n_{j=1}\frac{\beta_j}{\zeta_0-t_j}. 


$$ 


 


Пусть $\wt\eta_0$ настолько мало, что 


$\wt\eta_0<\frac{1-\Tilde\beta}{C+(n+1)/(\sqrt{2}\varepsilon)}$. 


Тогда $C+\frac{n-1}{\sqrt{2}\varepsilon}<


\frac{1-\Tilde\beta}{\Tilde\eta_0}$ 


и 


$$


|F_1(\zeta_0)|=\bigg|\frac{N(\zeta_0)}{M(\zeta_0)}-\sum^n_{j=1} 


\frac{\beta_j}{\zeta_0-t_j} \bigg|\le C+\bigg|\sum_{j\ne k} 


\frac{\beta_j}{\zeta_0-t_j} \bigg|+\bigg|\frac{\beta_k}{\zeta_0-t_k} 


\bigg| \le 


C+\frac{n-1}{\sqrt{2}\omega}+\frac{\wt\beta}{\wt\eta_0} 


<\frac{1}{\wt\eta_0}. 


$$


 


б) Пусть теперь $\xi_0$ не попадает ни в одну из 


$\varepsilon$-окрестностей точек $t_j$, 


$j=\ov{1,n}$. Тогда 


$$ 


\bigg|\sum^n_{j=1}\frac{\beta_j}{\zeta_0-t_j} \bigg| 


\le\frac{n}{\varepsilon}. 


$$ 


Если $\wt\eta_0<(C+n/\omega)^{-1}$, то 


$$ 


|F_1(\zeta_0)|\le C+n/\varepsilon <1/\wt\eta_0. 


$$ 


 


Таким образом, если $\wt\eta<\min\big[\frac{1-\Tilde\beta} 


{C+(n-1)/(\sqrt{2}\varepsilon)},\,(C+n/\varepsilon)^{-1}\big]$, то 


$|F_1(\zeta_0)|<1/\wt\eta_0$ 


как в случае а), так и в случае б). Из последнего 


неравенства следует, что значения функции 


$$ 


F(\zeta_0)=\frac{N(\zeta_0)}{M(\zeta_0)}+\frac{1}{i\eta_0} 


-\sum^n_{j=1}\frac{\beta_j}{\zeta_0-t_j} 


=-\frac{i}{\eta_0}+F_1(\zeta_0) 


$$ 


при $\zeta_0\in T_1$ лежат в нижней полуплоскости, следовательно, 


\begin{equation} 


\bigg|\int_{T_1}d\arg F(\zeta_0) \bigg|\le\pi. 


\end{equation} 


 


2. Пусть $\zeta_0=R e^{i\varphi}\in T_2$. Величину $R$ пока не 


фиксируем и выберем позже. Тогда 


$$ 


\bigg|\frac{\beta_j}{\zeta_0-t_j}-\frac{\beta_j}{\zeta_0} \bigg|= 


\bigg|\frac{\beta_j t_j}{R e^{i\varphi} 


(R e^{i\varphi}-t_j)} \bigg|\le 


\frac{|\beta_j|\,|t_j|}{R(R-1)}<\frac{1}{(R-1)^2}. 


$$ 


 


Учитывая, что $\sum\limits^n_{j=1}\beta_j=3$ и 


$$ 


N(\zeta_0)/M(\zeta_0)=-1/\zeta_0+O(R^{-2}),\quad 


R=|\zeta_0|\to\infty, 


$$ 


равномерно по $\varphi=\arg\zeta_0$, можем записать


$F(\zeta_0)=F_2(\zeta_0)+O(R^{-2})$, где 


$$ 


F_2(\zeta_0)=\frac{1}{i\eta_0}-\frac{2}{\zeta_0}= 


\frac{\zeta_0-2i\eta_0}{i\eta_0\zeta_0}= 


\frac{\ov\zeta_0}{i\eta_0\zeta_0}=- 


\frac{i\ov\zeta{}^2_0}{\eta_0R^2}, 


$$ 


откуда 


$$ 


\arg F_2(\zeta_0)=-2\varphi-\pi/2,\quad 


|F_2(\zeta_0)|\ge R^{-1}. 


$$ 


 


Таким образом, $F(\zeta_0)=F_2(\zeta_0)+o(F_2(\zeta_0))$, 


$R=|\zeta_0|\to\infty$, равномерно 


по $\varphi=\arg\zeta_0$. Так как $\int\limits^{2\pi}_0 d\arg F_2(R 


e^{i\varphi})=2\pi$, 


то при достаточно больших $R$ и малых $\wt\eta_0$ 


\begin{equation} 


\int_{T_2}d\arg F(\zeta_0)>\pi. 


\end{equation} 


Из (7) и (8) следует, что $\int\limits_{\partial Q}d\arg 


F(\zeta_0)\ne0$. 


 


Таким образом, вращение векторного поля, определяемого левой частью 


равенства (6), на $\partial Q$ отлично от нуля. Следовательно, это 


векторное поле 


в некоторой точке $\zeta_0\in Q$ обращается в нуль. Это доказывает 


разрешимость 


уравнения Гахова (6) в $D_\zeta$. 


 


 


\section{Структура множества корней уравнения Гахова} 


 


Рассмотрим уравнение Гахова 


\begin{equation} 


\frac{-2}{\zeta_0-\ov\zeta_0} 


\int^1_{-1}\frac{h(t)|t-\zeta_0|^4}{\Pi(t)(t-\zeta_0)}dt+ 


\sum^n_{j=1}\frac{\beta_j}{\zeta_0-t_j} 


\int^1_{-1}\frac{h(t)|t-\zeta_0|^4}{\Pi(t)(t-\zeta_0)}dt+ 


\int^1_{-1}\frac{h(t)|t-\zeta_0|^4}{\Pi(t)(t-\zeta_0)^2}dt=0. 


\end{equation} 


 


Введем обозначение $h(t)/\Pi(t)=G(t)$. 


Принимая во внимание равенства


\begin{gather*} 


\frac{|t-\zeta_0|^4}{(t-\zeta_0)}= 


(t-\zeta_0)(t-\ov\zeta_0)^2 (t-\zeta_0)= 


t^3+t^2(-\zeta_0-2\ov\zeta_0)+ 


t(\ov\zeta{}^2_0+2\ov\zeta_0\zeta_0)-\zeta_0\ov\zeta{}^2_0,\\ 


% 


\frac{|t-\zeta_0|^4}{(t-\zeta_0)^2}= 


t^2-2\ov\zeta_0t+\ov\zeta{}^2_0, 


\end{gather*} 


перепишем уравнение (9) в следующем виде: 


\begin{multline*} 


\frac{-2}{\zeta_0\ov\zeta_0}\bigg\{\int^1_{-1}G(t)t^3dt- 


(\zeta_0+2\ov\zeta_0)\int^1_{-1}G(t)t^2dt+ 


(\ov\zeta{}^2_0+2\ov\zeta_0\zeta_0)\int^1_{-1}G(t)t\,dt - \\ 


% 


-\zeta_0\ov\zeta{}^2_0\int^1_{-1}G(t)dt\bigg\}+ 


\sum^n_{j=1}\frac{\beta_j}{\zeta_0-t_j}\bigg\{\int^1_{-1} 


G(t)t^3dt-(\zeta_0+2\ov\zeta_0)\int^1_{-1}G(t)t^2dt+ \\ 


% 


+(\ov\zeta{}^2_0+2\ov\zeta_0\zeta_0)\int^1_{-1}G(t)t\,dt- 


\zeta_0\ov\zeta{}^2_0\int^1_{-1}G(t)dt\bigg\}+ \\ 


% 


+\int^1_{-1}G(t)t^2dt-2\ov\zeta_0\int^1_{-1}G(t)t\,dt 


+\ov\zeta{}^2_0\int^1_{-1}G(t)dt=0. 


\end{multline*} 


 


Обoзнaчaя $\int\limits^1_{-1}G(t)t^k dt=a_k$, $k=0,1,2,3$, 


а левую часть уравнения (9) через $F(\zeta_0)$, и приводя подобные 


слагаемые, получим 


\begin{multline*} 


F(\zeta_0)=\ov\zeta{}^3_0\bigg[\,\sum^n_{j=1} 


\frac{\beta_j}{\zeta_0-t_j}(a_1-a_0\zeta_0)-a_0\bigg]+ 


\ov\zeta{}^2_0\bigg[3a_0\zeta_0-a_0\zeta^2_0-2a_2 


\sum^n_{j=1}\frac{\beta_j}{\zeta_0-t_j}- \\ 


% 


-a_1\sum^n_{j=1}\frac{\beta_j}{\zeta_0-t_j}\zeta_0\bigg] 


+\ov\zeta_0\bigg[-5a_2-6a_1\zeta_0+3a_2\sum^n_{j=1} 


\frac{\beta_j}{\zeta_0-t_j}+2a_1\sum^n_{j=1} 


\frac{\beta_j}{\zeta_0-t_j}\zeta^2_0- \\ 


% 


-a_3\sum^n_{j=1}\frac{\beta_j}{\zeta_0-t_j}\bigg] 


-2a_3+3a_2\zeta_0+a_3\sum^n_{j=1}\frac{\beta_j}{\zeta_0-t_j} 


\zeta_0-a_2\sum^n_{j=1}\frac{\beta_j}{\zeta_0-t_j}\zeta^2_0. 


\end{multline*} 


Отсюда видно, что $F(\zeta_0)$ --- 4-аналитическая функция в верхней 


полуплоскости $D_\zeta$. Из доказанного в предыдущем параграфе следует, 


что 


множество $E$ нулей функции $F(\zeta_0)$ в $D_\zeta$ непусто. 


 


\begin{lemnonum} 


Множество $E$ не имеет предельных точек на границе $D_\zeta$. 


\end{lemnonum} 


 


{\bf Доказательство.}


Рассмотрим поведение $F(\zeta_0)$ на $\partial D_\zeta$. Повторяя 


рассуждения п.\,1) из 


предыдущего параграфа, получим, что при 


\begin{equation} 


\im\zeta_0=\eta_0<\min\bigg\{\frac{1-\wt\beta} 


{C(n-1)/(2\varepsilon)},\; \frac{1}{n/(2\varepsilon)+C}\bigg\} 


\end{equation} 


значения $F(\zeta_0)$ лежат в нижней полуплоскости, т.\,е. 


$|F(\zeta_0)|\ne0$ при 


выполнении условия (10). 


 


Пусть $R=|\zeta_0|\to\infty$. Тогда 


$F(\zeta_0)=F_2(\zeta_0)+o(F_2(\zeta_0))$, $R\to\infty$, 


равномерно по $\arg\zeta_0$, где 


$F_2(\zeta_0)=\frac{1}{i\eta_0}-\frac{2}{\zeta_0}$, 


$|F_2(R e^{i\varphi})|>R^{-1}$. Из последнего неравенства видно, что 


$|F_2(\zeta_0)|\ne0$, а 


следовательно, и $|F(\zeta_0)|\ne0$ при достаточно больших 


$R$.$\quad\square$ 


 


Воспользуемся одним результатом из теории полианалитических 


функций.


\medskip 


 


{\bf Теорема A.} ([7]).


{\it Пусть $f(z)$ --- $n$-аналитическая в некоторой области 


$D$ функция, отличная от тождественного нуля, и $E$ --- множество ее 


нулей в $D$. Если $a$ --- неизолированная точка из $E$, то в достаточно 


малой окрестности $\delta$ этой точки множество $E\setminus\{a\}$ 


является 


объединением не более чем $2n-2$ аналитических дуг.} 


\medskip 


 


Используя этот результат, получим следующее утверждение. 


 


\begin{thmnonum} 


Множество $E$ корней уравнения Гахова $(6)$ в $D_\zeta$ непусто, 


компактно и не имеет предельных точек на границе верхней полуплоскости. 


Это множество можно представить в виде $E=E_1\cup E_2$, 


$E_1\cap E_2\ne\emptyset$, где $E_1$ состоит из не более чем конечного 


числа точек, $E_2$ 


является объединением не более чем конечного числа попарно 


непересекающихся жордановых аналитических дуг. 


 


Если $a$ --- неизолированная точка множества $E$, то в достаточно 


малой окрестности этой точки множество $E\setminus\{a\}$ является 


объединением не более чем шести аналитических дуг. 


\end{thmnonum} 


 


{\bf Доказательство.} Функция $F(\zeta_0)$ в уравнении Гахова 


непрерывна, следовательно, $E$ -- замкнутое множество в $D_\zeta$. Так 


как $E$ не 


имеет предельных точек на границе верхней полуплоскости, то существует 


компакт $K$ такой, что $E\in K$. Значит, $E$ ограничено. Из 


ограниченности 


и замкнутости $E$ следует, что $E$ компактно в $\ov{\Bbb C}$. Для любой 


точки 


$a\in E$ по теореме A существует окрестность $\delta$ такая, что 


$E\cap\delta$ состоит 


либо из единственной точки $\{a\}$, либо из объединения не более чем 


шести аналитических дуг.$\quad\square$


 


 


\begin{thebibliography}{9} 


 


\bibitem{1} 


Монахов В.Н. {\em Краевые задачи со свободными границами для 


эллиптических систем уравнений}. -- Новосибирск: СО АН СССР, 1977. -- 


424 с. 


\bibitem{2} 


Насыров С.Р. {\em Смешанная обратная краевая задача на римановых 


поверхностях} // Изв. вузов. Математика. -- 1990. -- \No\,10. -- 


С.\,25--36. 


\bibitem{3} 


Авхадиев Ф.Г. {\em Конформные отображения и краевые задачи}. -- 


Казань: Казанский фонд ``Математика'', 1996. -- 216 с. 


\bibitem{4} 


Гахов Ф.Д. {\em Краевые задачи}. -- М.: Наука, 1977. -- 640 с. 


\bibitem{5} 


Аксентьев Л.А. {\em Связь внешней обратной краевой задачи с 


внутренним радиусом} // Изв. вузов. Математика. -- 1984. 


-- \No\,2. -- С.\,3--11. 


\bibitem{6} 


Киселев А.В., Насыров С.Р. {\em О структуре множества корней 


уравнения Ф.Д.\,Гахова для односвязной и многосвязной областей} // 


Труды семин. по краев. задачам. -- Казань: Изд-во Казанск. ун-та, 1990. 


-- 


Вып.\,24. -- С. 105--115. 


\bibitem{7} 


Balk М.В. {\em Polyanalytic functions}. -- 


Berlin: Akad. Verl., 1991. -- 197 р. 


 


\end{thebibliography} 


 


\vskip 2cm 


 


\post{\it Казанский государственный университет}{\it Поступила} 


\post{}{17.06.2002} 


 


\label{end} 


\end{document} 


